% VARIANT of sections/07-vector.tex (frame 1) — L: κάθετη ροή (stacked bullets).
% Υπόθεση: αντικαθιστά το κάτω two-column split με ενιαία κάθετη ροή
% (πρώτα η σημείωση VLMAX, μετά η "Επίδραση στον πυρήνα") και προσθέτει
% οπτικό separator ανάμεσα στο Tier 2 (setup) και στο Tier 3 (λύση στο leaf).
\begin{frame}{Tier 2+3: Διανυσματοποίηση βρόχων ανασύνθεσης και leaf}
  \pdfpcnote{budget ≈75s}
  % --- Tier 2: τι έμεινε ---
  \begin{itemize}
    \item \textbf{Tier 2:} διανυσματοποιεί το \BigO{n} (ανασύνθεση), όχι το \BigO{n^2} leaf
  \end{itemize}
  \vspace{0.3em}
  {\color{gray!45}\hrule}
  \vspace{0.6em}
  % --- Tier 3: η λύση στο leaf ---
  \HL{Tier 3: \vclmul{} υπολογίζει πολλαπλά $64\times64$ carry-less γινόμενα.}
  \vspace{0.3em}
  \begin{center}
    \resizebox{0.82\linewidth}{!}{% Fig: scalar Zbc leaf vs vector Zvbc (vclmul.vx) leaf.
% Presentation-oriented companion to schoolbook-vs-clmul.tex, used by
% presentation-v1/sections/07-vector (Tier 1/2 scalar leaf vs Tier 3 vector leaf).
% Semantic colour: left = costly scalar path (figthreat), right = vector
% solution (figsolution); transition arrow = critical path (figedge/accent).
\tikzsetnextfilename{ch6-polymul-leaf-scalar-vs-vclmul}
\begin{tikzpicture}[
  box/.style={draw, rounded corners=2pt, minimum width=5.0cm,
    minimum height=2.2cm, align=left, figT/mono, figLW/base},
  caption/.style={figT/annot, below=2pt},
]

% Left box: scalar Zbc leaf (costly path)
\node[box, fill=figthreatfill, draw=figthreat] (scalar) at (0,0) {%
  \begin{tabular}{@{}l@{}}
  for i in 0..n-1:\\
  \quad for j in 0..n-1:\\
  \quad\quad lo = clmul(a[i], b[j])\\
  \quad\quad hi = clmulh(a[i], b[j])\\
  \quad\quad r[i+j] \^{}= lo\\
  \quad\quad r[i+j+1] \^{}= hi
  \end{tabular}
};
\node[figT/title, above] at (scalar.north) {Scalar \textsf{Zbc} leaf};
\node[caption] at (scalar.south) {$n^2$ scalar clmul/clmulh};

% Right box: vector Zvbc vclmul.vx leaf (solution)
\node[box, fill=figsolutionfill, draw=figsolution] (vector) at (8.0,0) {%
  \begin{tabular}{@{}l@{}}
  for i in 0..n-1:\\
  \quad vb  = b[0..n-1]\\
  \quad vlo = vclmul.vx(vb, a[i])\\
  \quad vhi = vclmulh.vx(vb, a[i])\\
  \quad vr \^{}= combine(vlo, vhi)
  \end{tabular}
};
\node[figT/title, above] at (vector.north) {Vector \textsf{Zvbc} leaf};
\node[caption] at (vector.south) {1 vclmul.vx/vclmulh.vx per $i$ (VLMAX $=T$)};

% Transition arrow (critical path): inner j-loop collapses to one vector op
\draw[figedge/accent]
  (scalar.east) -- node[above, figT/label] {$j$-loop $n\to1$}
  (vector.west);

\end{tikzpicture}
}
  \end{center}
  % src: figures/ch6-polymul/leaf-scalar-vs-vclmul.tex
  \vspace{0.3em}
  \pause
  \begin{itemize}
    \item Σε \VLEN{}=256, \code{LMUL=m4}: \HL{VLMAX $= T = 16$}, αρκεί ένα vector iteration ανά $i$
    \item \textbf{Αποτέλεσμα:} \HL{\speedup{264}} στον πυρήνα πολ/σμού (\HQC-128, Tier 3)
  \end{itemize}
  \pdfpcnote{Tier 2: ο κρίσιμος βρόχος ανασύνθεσης περνά από scalar XOR σε vxor.vv, γραμμικό κόστος. Ο leaf clmul/clmulh (O(n^2)) ήταν ήδη ο κυρίαρχος όρος, όχι κάτι που εμφανίζεται μετά το Tier 2.}
  \pdfpcnote{.vx μορφή: για κάθε i, το scalar a[i] πολλαπλασιάζεται carry-less με ολόκληρο τον πίνακα b[0..n-1]. Σε VLMAX=16=T ο inner j-loop γίνεται μία iteration, μηδενίζοντας τον scalar inner loop.}
  \pdfpcnote{264 = polymul HQC-128 Tier 3 (σωρευτικά vs scalar_ref). Το HQC-256 (x436) πάει στο results slide.}
\end{frame}

\begin{frame}{Μετά την βελτιστοποίηση: \HL{\en{Keccak}} γίνεται το νέο \en{Amdahl} floor (\pct{49.2})}
  \pdfpcnote{budget ≈45s}
  \HL{Ο πολ/σμός (αρχικά ${\approx}$\pct{97}) υποχωρεί σε ${\approx}$\pct{33} (Tier 1) $\to$ ${\approx}$\pct{22} (Tier 3). Κυρίαρχο πλέον το \en{Keccak}.}
  \vspace{0.3em}
  \begin{columns}[T,onlytextwidth]
    \begin{column}{0.5\textwidth}
      \centering
      \resizebox{0.9\linewidth}{!}{\tikzsetnextfilename{ch6-polymul-pie-scalar-zb}
\begin{tikzpicture}[
    every node/.style={font=\sffamily},
]

% --- Palette via canonical component tokens (cross-figure consistency) ---
\colorlet{colMul}{figkPolymul}
\colorlet{colKeccak}{figkKeccak}
\colorlet{colDecode}{figkDecode}
\colorlet{colMem}{figkMisc}
\colorlet{textmain}{black!75}
\colorlet{textsub}{figneutral}

% --- Parameters ---
\def\mainCX{0}
\def\mainCY{0}
\def\outerR{2.8}
\def\innerR{1.7}

% Helper: ring sector (cx, cy, innerR, outerR, startAngle, endAngle, color)
\newcommand{\ringsectorB}[7]{%
    \fill[#7, draw=white, line width=1.2pt]
        ([shift={(#1,#2)}] #5:#3) arc (#5:#6:#3)
        -- ([shift={(#1,#2)}] #6:#4) arc (#6:#5:#4) -- cycle;
}

% MAIN DONUT — scalar_zb (Tier 1) L128 profile (keccak=ref, decode=ref)
% Polymul = 36.8%  → 132.5°
% Keccak  = 26.5%  →  95.4°
% Decode  = 16.2%  →  58.3°
% Other   = 20.5%  →  73.8°
%
% Angles: start=90, clockwise (90 -> -270)
% Polymul:  90.0  -> -42.5
% Keccak: -42.5  -> -137.9
% Decode: -137.9 -> -196.2
% Other:  -196.2 -> -270.0

\ringsectorB{\mainCX}{\mainCY}{\innerR}{\outerR}{90}{-42.5}{colMul}
\ringsectorB{\mainCX}{\mainCY}{\innerR}{\outerR}{-42.5}{-137.9}{colKeccak}
\ringsectorB{\mainCX}{\mainCY}{\innerR}{\outerR}{-137.9}{-196.2}{colDecode}
\ringsectorB{\mainCX}{\mainCY}{\innerR}{\outerR}{-196.2}{-270}{colMem}

% Labels inside sectors
\pgfmathsetmacro{\labelR}{(\innerR + \outerR)/2}

% Polymul label (mid-angle = (90 + (-42.5))/2 = 23.7)
\node[font=\sffamily\bfseries\small, text=white] at
    ([shift={(\mainCX,\mainCY)}] 23.7:\labelR) {36.8\%};

% Keccak label (mid-angle = (-42.5 + (-137.9))/2 = -90.2)
\node[font=\sffamily\bfseries\small, text=white] at
    ([shift={(\mainCX,\mainCY)}] -90.2:\labelR) {26.5\%};

% Decode label (mid-angle = (-137.9 + (-196.2))/2 = -167.1)
\node[font=\sffamily\bfseries\small, text=white] at
    ([shift={(\mainCX,\mainCY)}] -167.1:\labelR) {16.2\%};

% Other label (mid-angle = (-196.2 + (-270))/2 = -233.1)
\node[font=\sffamily\bfseries\small, text=white] at
    ([shift={(\mainCX,\mainCY)}] -233.1:\labelR) {20.5\%};

% Center text
\node[font=\sffamily\bfseries\large, text=textmain] at (\mainCX, \mainCY+0.15)
    {HQC-128};
\node[font=\sffamily\small, text=textsub] at (\mainCX, \mainCY-0.25)
    {Scalar Zb (Tier~1)};

% LEGEND (bottom)
\def\legY{-\outerR-0.9}
\def\legBoxSize{0.25}

\fill[colMul] (-3.0, \legY) rectangle +(\legBoxSize, \legBoxSize);
\node[font=\sffamily\scriptsize, text=textmain, anchor=west] at (-3.0+\legBoxSize+0.1, \legY+0.12)
    {Polymul};

\fill[colKeccak] (-1.2, \legY) rectangle +(\legBoxSize, \legBoxSize);
\node[font=\sffamily\scriptsize, text=textmain, anchor=west] at (-1.2+\legBoxSize+0.1, \legY+0.12)
    {Keccak/Hashing};

\fill[colDecode] (1.2, \legY) rectangle +(\legBoxSize, \legBoxSize);
\node[font=\sffamily\scriptsize, text=textmain, anchor=west] at (1.2+\legBoxSize+0.1, \legY+0.12)
    {Decode};

\fill[colMem] (3.0, \legY) rectangle +(\legBoxSize, \legBoxSize);
\node[font=\sffamily\scriptsize, text=textmain, anchor=west] at (3.0+\legBoxSize+0.1, \legY+0.12)
    {Other};

\end{tikzpicture}
}
    \end{column}
    \begin{column}{0.5\textwidth}
      \centering
      \resizebox{0.9\linewidth}{!}{\tikzsetnextfilename{ch6-polymul-pie-vector-zvk}
\begin{tikzpicture}[
    every node/.style={font=\sffamily},
]

% --- Palette via canonical component tokens (cross-figure consistency) ---
\colorlet{colMul}{figkPolymul}
\colorlet{colKeccak}{figkKeccak}
\colorlet{colDecode}{figkDecode}
\colorlet{colMem}{figkMisc}
\colorlet{textmain}{black!75}
\colorlet{textsub}{figneutral}

% --- Parameters ---
\def\mainCX{0}
\def\mainCY{0}
\def\outerR{2.8}
\def\innerR{1.7}

% Helper: ring sector (cx, cy, innerR, outerR, startAngle, endAngle, color)
\newcommand{\ringsectorC}[7]{%
    \fill[#7, draw=white, line width=1.2pt]
        ([shift={(#1,#2)}] #5:#3) arc (#5:#6:#3)
        -- ([shift={(#1,#2)}] #6:#4) arc (#6:#5:#4) -- cycle;
}

% MAIN DONUT — vector_zvk (Tier 3) L128 profile (keccak=ref, decode=ref)
% Polymul = 21.8%  →  78.5°
% Keccak  = 49.2%  → 177.1°
% Decode  = 24.0%  →  86.4°
% Other   =  5.1%  →  18.4°
%
% Angles: start=90, clockwise (90 -> -270)
% Polymul:  90.0  ->  11.5
% Keccak:  11.5  -> -165.6
% Decode: -165.6  -> -252.0
% Other:  -252.0  -> -270.0

\ringsectorC{\mainCX}{\mainCY}{\innerR}{\outerR}{90}{11.5}{colMul}
\ringsectorC{\mainCX}{\mainCY}{\innerR}{\outerR}{11.5}{-165.6}{colKeccak}
\ringsectorC{\mainCX}{\mainCY}{\innerR}{\outerR}{-165.6}{-252.0}{colDecode}
\ringsectorC{\mainCX}{\mainCY}{\innerR}{\outerR}{-252.0}{-270}{colMem}

% Labels inside sectors
\pgfmathsetmacro{\labelR}{(\innerR + \outerR)/2}

% Polymul label (mid-angle = (90 + 11.5)/2 = 50.8)
\node[font=\sffamily\bfseries\small, text=white] at
    ([shift={(\mainCX,\mainCY)}] 50.8:\labelR) {21.8\%};

% Keccak label (mid-angle = (11.5 + (-165.6))/2 = -77.1)
\node[font=\sffamily\bfseries\small, text=white] at
    ([shift={(\mainCX,\mainCY)}] -77.1:\labelR) {49.2\%};

% Decode label (mid-angle = (-165.6 + (-252.0))/2 = -208.8)
\node[font=\sffamily\bfseries\small, text=white] at
    ([shift={(\mainCX,\mainCY)}] -208.8:\labelR) {24.0\%};

% Other label (mid-angle = (-252.0 + (-270))/2 = -261.0)
\node[font=\sffamily\bfseries\footnotesize, text=white] at
    ([shift={(\mainCX,\mainCY)}] -261.0:\labelR) {5.1\%};

% Center text
\node[font=\sffamily\bfseries\large, text=textmain] at (\mainCX, \mainCY+0.15)
    {HQC-128};
\node[font=\sffamily\small, text=textsub] at (\mainCX, \mainCY-0.25)
    {Vector Zvk (Tier~3)};

% LEGEND (bottom)
\def\legY{-\outerR-0.9}
\def\legBoxSize{0.25}

\fill[colMul] (-3.0, \legY) rectangle +(\legBoxSize, \legBoxSize);
\node[font=\sffamily\scriptsize, text=textmain, anchor=west] at (-3.0+\legBoxSize+0.1, \legY+0.12)
    {Polymul};

\fill[colKeccak] (-1.2, \legY) rectangle +(\legBoxSize, \legBoxSize);
\node[font=\sffamily\scriptsize, text=textmain, anchor=west] at (-1.2+\legBoxSize+0.1, \legY+0.12)
    {Keccak/Hashing};

\fill[colDecode] (1.2, \legY) rectangle +(\legBoxSize, \legBoxSize);
\node[font=\sffamily\scriptsize, text=textmain, anchor=west] at (1.2+\legBoxSize+0.1, \legY+0.12)
    {Decode};

\fill[colMem] (3.0, \legY) rectangle +(\legBoxSize, \legBoxSize);
\node[font=\sffamily\scriptsize, text=textmain, anchor=west] at (3.0+\legBoxSize+0.1, \legY+0.12)
    {Other};

\end{tikzpicture}
}
    \end{column}
  \end{columns}
  % src: figures/ch6-polymul/{pie-scalar-zb,pie-vector-zvk}.tex (HQC-128 instret profile)
  \pdfpcnote{Scalar Zb (Tier 1): polymul 32.7, Keccak 33.7, decode 15.3, other 18.2. Vector Zvk (Tier 3): polymul 21.8, Keccak 49.2, decode 24.0, other 5.1. Αυτό αιτιολογεί γιατί τα επόμενα βήματα στοχεύουν Keccak/decode.}
\end{frame}
